\begin{figure}[t]
\centering
\begin{tikzpicture}[x=1cm,y=1cm,>=Latex,font=\small]
  \draw[->,thick] (0,0)--(13,0) node[right] {\(s\)};
  \foreach \x/\y in {0.8/2.6,3.8/5.6,7.3/9.1} {
    \fill[blue!9] (\x,0.15) rectangle (\y,0.82);
    \draw[blue!60!black] (\x,0.15) rectangle (\y,0.82);
    \node at ({(\x+\y)/2},0.49) {test: \(\ell_n\)};
    \draw (\x,0.07)--(\x,-0.07);
    \draw (\y,0.07)--(\y,-0.07);
  }
  \node[below] at (0.8,-0.1) {\(\sigma_1\)};
  \node[below] at (2.6,-0.1) {\(\sigma_1+\ell_n\)};
  \node[below] at (3.8,-0.1) {\(\sigma_2\)};
  \node[below] at (5.6,-0.1) {\(\sigma_2+\ell_n\)};
  \node[below] at (7.3,-0.1) {\(\sigma_J\)};
  \node[below] at (9.1,-0.1) {\(\sigma_J+\ell_n\)};
  \node at (6.45,0.48) {\(\cdots\)};
  \draw[<->] (2.6,1.1)--(3.8,1.1)
    node[midway,above] {\(n\)};
  \draw[<->] (9.1,1.1)--(12.2,1.1)
    node[midway,above] {\(n+\frac{1}{2a}\log\frac{100}{8}\)};
  \draw[dashed] (12.2,-0.07)--(12.2,0.85);
  \node[below] at (12.2,-0.1) {\(\sigma_{\mathrm f}\)};
  \node[align=center,font=\footnotesize] at (6.5,-1.15)
    {All failures are known at \(T_J\). The final gap controls\\
     the inverse-angle factor in the change back to the original law.};
\end{tikzpicture}
\caption{One batch at a fixed grid center, in inner time (schematic).
There are \(J=I_n-1\) mass tests, each of length \(\ell_n\),
with gaps \(n\) between consecutive tests. The coordinates are
\(s_z(S_i)=\sigma_i\), \(s_z(T_i)=\sigma_i+\ell_n\), and
\(s_z(\tau_{\mathrm f})=\sigma_{\mathrm f}\); the final coordinate
corresponds to conformal radius \(8r_{n,I_n}\).
The first restart is preceded by a mixing gap as well.
The tests are conditioned on their actual past; independence is unnecessary.}
\label{fig:batch-timeline}
\end{figure}
